% Marco teórico, parte 2 (pessoa 5): velocity, produtividade e conversão de estimativa em custo.
% A parte 1 (histórias de usuário, story points e Planning Poker) fica em 02a-historias-pontos.tex.
% Os valores numéricos desta parte são hipotéticos e seguem o exemplo da Equação (eq:previsao).
% Convenção desta parte: sem travessões e sem ponto e vírgula. Citações com vários autores
% usam \citeonline, porque \cite separa os sobrenomes com ponto e vírgula no padrão ABNT.

\subsection{\textit{Velocity}}
\label{sub:velocity}

Saber o tamanho de uma história não diz, sozinho, quando ela vai ficar pronta.
Quem responde a essa pergunta é a \textit{velocity}, isto é, a quantidade de
trabalho que a equipe consegue concluir em uma iteração. Quando as histórias são
estimadas em \textit{story points}, a \textit{velocity} é simplesmente a soma dos
pontos das histórias terminadas no \textit{sprint} \cite{cohn2005}. A regra de
contagem é rígida de propósito. Uma história que ficou 90\% pronta vale zero
pontos naquele \textit{sprint}, porque ainda vai consumir esforço no seguinte, e
contá-la pela metade daria a impressão de um avanço que não aconteceu
\cite{cohn2005}.

Há um efeito curioso nessa medida. Se a equipe subestima as histórias de forma
sistemática, ela também termina menos pontos por iteração do que esperava, e a
divisão entre as duas grandezas continua dando um prazo coerente.
\citeonline{cohn2005} vê nisso uma forma de autocorreção da estimativa. Mas a
compensação depende de duas condições que nem sempre se mantêm. A primeira é que
a equipe use a escala de pontos sempre do mesmo jeito. A segunda é que a própria
equipe não mude muito. Basta trocar parte dos integrantes, adotar uma tecnologia
nova ou rever o critério de pontuação para que o histórico deixe de servir de
referência.

Na prática, a \textit{velocity} oscila de um \textit{sprint} para outro. Férias,
feriados, incidentes em produção e histórias que não cabem na iteração bastam
para tirá-la do lugar. Por isso, \citeonline{cohn2005} sugere trabalhar com a
média das últimas iterações e comunicar prazos como uma faixa, e não como um
número fechado. Para uma equipe sem histórico, o autor aponta três caminhos.
Ela pode usar dados de um projeto parecido, rodar algumas iterações antes de
prometer qualquer prazo ou estimar a capacidade disponível quebrando as
histórias em tarefas medidas em horas. Se $v_1, \dots, v_k$ são as
\textit{velocities} das $k$ últimas iterações, a média e os extremos observados
são
%
\begin{equation}
  \bar{v} = \frac{1}{k}\sum_{i=1}^{k} v_i, \qquad
  v_{\min} = \min_i v_i, \qquad
  v_{\max} = \max_i v_i .
  \label{eq:velocity}
\end{equation}

Prever a \textit{velocity} virou, inclusive, tema de pesquisa.
\citeonline{hearty2009} representam a \textit{velocity} de projetos XP com uma
rede bayesiana dinâmica, que se recalibra a cada iteração e, assim, melhora a
previsão à medida que o projeto gera dados. \citeonline{choetkiertikul2018}
seguem outro caminho e usam informações de iterações e de \textit{issues} de
projetos de código aberto para prever, com o \textit{sprint} ainda em andamento,
se a equipe vai entregar o que planejou. Os dois trabalhos partem da mesma
constatação. Tratar a \textit{velocity} como um número fixo esconde boa parte do
risco da previsão.

\subsection{Produtividade e capacidade da equipe}
\label{sub:produtividade}

Em economia, produtividade é a razão entre o que se produz e o que se gasta para
produzir. Na estimativa tradicional de software, ela costuma aparecer como
tamanho por esforço, em linhas de código por pessoa-mês, como nos modelos de
\citeonline{boehm1981}, ou em pontos de função por unidade de esforço, medida
que \citeonline{albrecht1983} relacionaram ao esforço de desenvolvimento. A
\textit{velocity} tem a mesma forma, pontos por iteração, mas carrega duas
diferenças que mudam a maneira de usá-la.

A primeira diferença está na unidade. Como vimos na
Seção~\ref{sub:storypoints}, o ponto de uma equipe não equivale ao ponto de
outra. Comparar a \textit{velocity} de duas equipes, portanto, é comparar
réguas diferentes, e \citeonline{cohn2005} desaconselha esse uso. O problema
piora quando a \textit{velocity} vira meta de desempenho. A equipe passa a
``entregar'' mais pontos sem entregar mais funcionalidade, bastando para isso
aumentar as estimativas, e a métrica perde justamente o que tinha de útil, que
era ajudar a prever. O estudo de caso de \citeonline{schull2023}, numa grande
montadora, mostra como é difícil levar indicadores de equipes ágeis aos níveis
de produto e portfólio sem distorcê-los. Já \citeonline{koch2008} mostra, em
projetos de código aberto, que a escolha da medida de esforço altera as
conclusões sobre a eficiência econômica do desenvolvimento.

A segunda diferença é que a \textit{velocity} mistura produtividade com
capacidade. Se ela cai, pode ser que a equipe tenha rendido menos, ou apenas que
tenha trabalhado menos dias. Uma forma simples de separar os dois efeitos é
corrigir a média pela capacidade prevista. Sendo $h_p$ as horas-pessoa
disponíveis no próximo \textit{sprint} e $\bar{h}$ a média de horas-pessoa das
iterações usadas no cálculo de $\bar{v}$, a \textit{velocity} esperada é
%
\begin{equation}
  v_{\text{esp}} = \bar{v}\,\frac{h_p}{\bar{h}} .
  \label{eq:capacidade}
\end{equation}
%
Uma equipe que costuma entregar 30 pontos e terá 20\% menos horas por causa de
um feriado e de uma licença deveria planejar algo perto de 24 pontos para o
\textit{sprint}. O ajuste supõe que a produtividade por hora não mude, o que é
razoável no curto prazo. Ela deixa de valer quando entra gente nova no time, porque as
horas disponíveis aumentam ao mesmo tempo em que o rendimento por hora cai,
enquanto os novos integrantes aprendem o sistema.

\subsection{Conversão da estimativa em prazo e custo}
\label{sub:custo}

A Equação~\ref{eq:previsao} já transforma o \textit{backlog} em número de
\textit{sprints} e em custo. Para que esse resultado sustente uma decisão de
investimento, ainda faltam duas coisas. É preciso saber de onde vem o custo da
equipe por \textit{sprint} e é preciso dizer o quanto a previsão pode errar.

O custo por \textit{sprint} vem do tempo das pessoas, que é o maior componente do
custo de um projeto de software \cite{boehm1981}. Com $n$ integrantes dedicados,
$h$ horas de trabalho de cada um no \textit{sprint} e um custo por hora $c_h$ que
já inclua salário, encargos e benefícios, tem-se
%
\begin{equation}
  c_s = n\,h\,c_h .
  \label{eq:custosprint}
\end{equation}
%
Os R\$\,60.000 por \textit{sprint} usados na Introdução correspondem, por
exemplo, a seis pessoas trabalhando 80 horas cada em um \textit{sprint} de duas
semanas, a R\$\,125 por hora. Gastos que não dependem do tamanho do time, como
licenças e infraestrutura, entram em $c_s$ quando crescem com a duração do
projeto, ou são somados uma única vez ao custo total quando não crescem.

Como a \textit{velocity} oscila, faz mais sentido prever uma faixa do que um
valor. Usando os extremos da Equação~\ref{eq:velocity} e arredondando o número de
\textit{sprints} para cima, já que ninguém paga meio \textit{sprint} de uma
equipe dedicada, chega-se a um cenário otimista e a um pessimista:
%
\begin{equation}
  N_{\min} = \left\lceil \frac{B}{v_{\max}} \right\rceil, \qquad
  N_{\max} = \left\lceil \frac{B}{v_{\min}} \right\rceil, \qquad
  N_{\min}\,c_s \le C \le N_{\max}\,c_s .
  \label{eq:intervalo}
\end{equation}
%
Suponha que as cinco últimas \textit{velocities} da equipe do exemplo tenham sido
26, 31, 28, 34 e 31 pontos. A média é exatamente 30, o mesmo valor usado na
Introdução, mas o \textit{backlog} de 240 pontos pode levar de 8 a 10
\textit{sprints}, e o custo fica entre R\$\,480.000 e R\$\,600.000. Os
R\$\,480.000 calculados antes são o melhor caso dessa faixa. Apresentar só esse
número a quem aprova o orçamento seria cair no excesso de confiança descrito por
\citeonline{jorgensen2005}. Vale lembrar que os extremos de cinco iterações não
são limites garantidos. Modelos que produzem uma distribuição de prazos, como a
rede bayesiana de \citeonline{hearty2009}, permitem dizer também qual é a
probabilidade de cada cenário.

A incerteza também muda ao longo do projeto. No começo, com requisitos ainda
vagos, ela é grande, e vai diminuindo conforme o produto é entendido e
construído. Esse padrão ficou conhecido como cone da incerteza \cite{boehm1981}.
\citeonline{little2006}, analisando dados de projetos reais, confirma que a
incerteza nas estimativas de prazo é alta, mas encontra um cone menos estreito
do que o modelo sugere e uma razão entre prazo real e prazo estimado com
distribuição assimétrica, de cauda longa para o lado dos atrasos. Para quem faz
a análise financeira, a lição é dupla. O cenário pessimista merece mais peso que
o otimista, e a previsão precisa ser refeita a cada \textit{sprint}, com o
\textit{backlog} e a \textit{velocity} atualizados.

\subsection{Acompanhamento do custo durante o projeto}
\label{sub:agileevm}

Depois que o orçamento é aprovado, as mesmas grandezas servem para acompanhar a
execução. \citeonline{sulaiman2006} adaptaram ao Scrum a gestão de valor
agregado (\textit{Earned Value Management}, EVM), técnica clássica de controle
de projetos, e chamaram o resultado de AgileEVM. A ideia central é medir o
progresso pelos pontos concluídos, sem depender de um cronograma detalhado de
tarefas. Considere $BAC$ o orçamento total (\textit{budget at completion}), $N$ o
número de \textit{sprints} planejados, $j$ os \textit{sprints} já concluídos,
$P_c$ os pontos entregues até agora e $AC$ o custo real gasto até o momento
(\textit{actual cost}). Com essas grandezas calculam-se o valor planejado $PV$,
o valor agregado $EV$ e os índices de desempenho de custo ($CPI$) e de prazo
($SPI$):
%
\begin{equation}
  PV = \frac{j}{N}\,BAC, \qquad
  EV = \frac{P_c}{B}\,BAC, \qquad
  CPI = \frac{EV}{AC}, \qquad
  SPI = \frac{EV}{PV} .
  \label{eq:agileevm}
\end{equation}

Voltando ao exemplo, com $BAC = \text{R\$}\,480.000$ e $N = 8$, imagine que, ao
fim do quarto \textit{sprint}, a equipe tenha entregado 112 dos 240 pontos e
gastado R\$\,250.000. Nesse caso, $PV = \text{R\$}\,240.000$,
$EV = \text{R\$}\,224.000$, $CPI \approx 0{,}90$ e $SPI \approx 0{,}93$. Em
outras palavras, o projeto entregou menos do que o planejado para a metade do caminho
e gastou mais do que vale o que entregou. Se esse ritmo continuar, o custo
estimado no término é $BAC/CPI \approx \text{R\$}\,536.000$. O valor cai dentro
da faixa da Equação~\ref{eq:intervalo}, mas fica bem acima dos R\$\,480.000 que
teriam sido prometidos se a previsão fosse pontual. O AgileEVM não resolve a
incerteza dos pontos, já que $EV$ depende de $B$, e $B$ muda sempre que o
\textit{backlog} muda. O que ele oferece é uma tradução. O progresso, que o time
enxerga em pontos, passa a ser descrito em custo e prazo, que é a linguagem de
quem decide sobre o investimento.

Olhando o caminho inteiro, da história de usuário até o orçamento, há três
conversões. A equipe estima o tamanho em pontos, a \textit{velocity}
transforma pontos em tempo e o custo do time transforma tempo em dinheiro. Cada
uma delas soma incerteza à anterior. Por isso, uma previsão financeira baseada
em \textit{story points} só ajuda de verdade quando chega a quem decide como
uma faixa de valores, revista a cada iteração, e não como um número único
tratado como certo.
